\pageheading{RESEARCH JUDGMENT}{What 1.0 closes, and what remains open}
The 1.0 release is a mathematical and software milestone. It closes selected
structural questions and provides a reproducible visual treatment; historical
priority remains a separately bounded literature question.

\begin{tabular}{@{}>{\raggedright\arraybackslash}p{0.46\linewidth}>{\raggedright\arraybackslash}p{0.47\linewidth}@{}}
\toprule
Closed in this release & Still open here\\\midrule
Positive lower and finite upper dilation type & Existence of an optimal type limit and logarithmic modulation\\
Joint analytic parameter dependence and compact tap bounds & General minimal-bank identifiability\\
Exact arithmetic parameter recovery up to modulus radical & Stable recovery within bounded prime classes\\
All-prime signed-atom detectors and positivity domain & Sharp sieve-root asymptotics and multiplicities\\
Global orbit-zeta continuation and shifted determinants & Functional equations and derivative-weighted zero sums\\
Integer coefficient and exact polyphase/tap machinery & Historical priority of the candidate extensions\\
\bottomrule\end{tabular}

\claim{The strongest direct function candidate}
The asymmetric dilation extension remains the strongest direct discovery
candidate: its parameters are intrinsic, its off-diagonal members cannot be
rescaled from the known diagonal, and its nonnegative real scale regime has now been
mapped. The arithmetic inverse and signed-resolution results add structural
depth. The orbit continuation is a worked application of established spectral
methods, with its exact seed and branch conventions verified locally.

\claim{A common finite-resolution lesson}
A weak tap can change asymptotic order while preserving compact convergence.
A large excluded prime can change an exact arithmetic parameter while hiding
within a small bounded-window error. A negative atom can evade low-degree
positivity and become detectable at a calculable larger degree.
These are three different proved mechanisms. Their shared resolution theme
does not make them identical constructions.

\pageheading{RELEASE VERIFICATION}{Evidence with declared ranges}
The independent audit reconstructed the five new proof chapters and checked
their final revisions. The release runner executed ten scripts once, with
25 pinned inputs and byte-identical isolated snapshots. All inputs remained
unchanged, every declared result was present, and the schema validator passed.

\[
\boxed{19{,}871\text{ exact assertion statements}\quad+
5{,}576\text{ numerical assertion statements passed}.}
\]
The exact total comprises 19,331 principal checks and 540 separately recorded
helper checks in the inverse reconstruction. Operational path and process
guards are excluded. Counts are assertion statements, not counts of theorems
or every scalar operation inside a grouped assertion.

The API contributes 8,038 exact assertions and 18 catalog rows through index
12. Dilation type and continuity add 6,091 exact checks. Arithmetic inverse
examples include rank 10,000, negative depth 500, and large excluded primes.
The multiple-prime run tests 102 positive parameter/depth cases, 1,395
orthogonality conditions, 381 recurrence cases, and 40 conditional pencils.
Spectral continuation tests classical unshifted and shifted anchors.

\claim{Recovered failures are retained}
The initial 0.1 environment and plotting issues, unsuitable absolute
determinant residual, and known one-dilation overlap remain documented.
In 1.0, two API validation defects were reproduced and fixed. A draft
sequence example had two miscopied integers and was corrected against the
exact catalog. Initial figure footer and label overlaps were repaired.
A direct publisher-source retrieval failed; an authenticated publisher
preview supplied the classical Laplace-uniqueness comparison instead.

Exact analytic tail envelopes are distinguished from floating-point rounding.
All decimal arithmetic and root approximations are ordinary high precision,
not interval certificates. No analytic Lean certificate is claimed.
The universal conclusions rely on the proofs, not on the number of checks.

\pageheading{RELEASE USE}{A reproducible research package}
The source archive includes the manuscript, all 18 vector figures and PNG
previews, figure generators, exact API, JSON/CSV catalog, inherited baseline
routes, the five new proof chapters, independent audit, and bounded release
evidence with input and output hashes.

The baseline route files preserve the original SF-DISCOVERY-001 claims and
open questions as historical evidence. The SF-DISCOVERY-100 chapters supersede
those questions where a new proof closes them. The current status is the
release matrix above, not an old route paragraph read in isolation.

\claim{Reproduce the work}
Use the versions listed in \texttt{requirements.txt}; execute
\texttt{checks/run\_release.py} in a fresh output directory. The checked
instructions give the bounded-run invocation and validator. Figure scripts
then rebuild the atlas from formulas and recorded numerical examples.
Two \texttt{pdflatex} passes compile the source.

The standard-library API is portable. Numerical checks additionally require
mpmath. No network access, external submission, or dependency installation is
part of the verification run. PDF byte identity can change with timestamps
and font environments; mathematical outputs retain their own provenance.

The remaining questions have executable starting points: type-limit analysis
for dilation, precision-controlled arithmetic inversion within bounded
parameter classes, sharper sieve-root asymptotics, and spectral weighted sums.
They remain research questions rather than hidden dependencies of the
completed release claims.
